\chapter{Métrica Principal de Validación}

Este capítulo define, justifica y reporta las métricas de validación del MVP.
La selección de indicadores responde a los objetivos SMART establecidos en el
Capítulo 3 y a los criterios de aceptación del Capítulo 6, asegurando que
los resultados sean comparables, reproducibles y significativos para los
tomadores de decisión en ProInversión.

\section{Marco de Métricas: Criterios de Selección}

Las métricas fueron seleccionadas bajo tres criterios:

\begin{enumerate}
    \item \textbf{Medibilidad objetiva:} cada indicador puede calcularse a
    partir de datos registrados por el sistema (logs, base de datos) o de
    un protocolo de evaluación manual replicable.
    \item \textbf{Alineación con el problema:} reflejan directamente los
    tres ejes identificados en el Capítulo 2: eficiencia operativa
    (tiempo), calidad del análisis (precisión/cobertura) y trazabilidad
    institucional.
    \item \textbf{Comparabilidad con el baseline:} permiten calcular la
    mejora respecto al proceso manual (320 horas, cobertura parcial,
    trazabilidad nula).
\end{enumerate}

Se distinguen dos categorías: una métrica principal de impacto operativo
y cinco métricas secundarias de calidad técnica.

\section{Métrica Principal: Tiempo de Screening Inicial}

\subsection{Definición formal}

El tiempo de screening se define como:
\[
T_{\text{screening}} = T_{\text{reporte}} - T_{\text{carga}}
\]
donde $T_{\text{carga}}$ es el timestamp del momento en que el PDF queda
almacenado en Cloud Storage, y $T_{\text{reporte}}$ es el timestamp de la
generación exitosa del reporte ejecutivo registrado en Cloud SQL.

\subsection{Justificación como métrica principal}

El tiempo de screening es el indicador de mayor impacto operativo porque:

\begin{itemize}
    \item Cuantifica directamente la reducción del esfuerzo humano.
    El proceso manual demanda $\approx 320$ horas (8 semanas de trabajo
    de un equipo multidisciplinario de 5 especialistas).
    \item Es interpretable por audiencias no técnicas: directivos,
    usuarios institucionales y auditores externos.
    \item Es sensible a mejoras o regresiones en cualquier componente
    del pipeline (segmentación, RAG, LLM, consolidación).
    \item Define si el sistema es viable para contextos reales de
    supervisión contractual, donde los plazos institucionales son
    estrictos.
\end{itemize}

\subsection{Objetivo SMART y resultado}

\begin{table}[H]
\centering
\caption{Métrica principal: tiempo de screening inicial}
\label{tab:metrica_principal}
\begin{tabular}{p{3.5cm} p{3.5cm} p{5.5cm}}
\toprule
\textbf{Parámetro} & \textbf{Valor} & \textbf{Observación} \\
\midrule
Baseline manual &
    320 horas &
    Proceso completo con equipo de 5 especialistas \\
Objetivo SMART &
    $\approx 1$ hora &
    Para contratos de 300+ páginas \\
Resultado Gemini 2.5 Pro &
    $\approx 1$ h 15 min &
    Contrato A (324 págs., 41 secciones) \\
Resultado Gemini 3.1 Preview &
    $\approx 2$ h 50 min &
    Mismo contrato; mayor latencia por modelo \\
Reducción lograda &
    $99.6\%$ &
    Respecto al baseline manual \\
Estado &
    \checkmark\ SUPERADO &
    El objetivo de $\approx 1$ hora se alcanzó con Gemini 2.5 Pro \\
\bottomrule
\end{tabular}
\end{table}

\subsection{Descomposición del tiempo por fase}

El tiempo total se distribuye entre las fases del pipeline. Los valores
corresponden al análisis del Contrato A con Gemini~2.5~Pro en
producción:

\begin{table}[H]
\centering
\caption{Distribución del tiempo de procesamiento por fase del pipeline}
\label{tab:tiempo_fases}
\begin{tabular}{p{1.5cm} p{4.5cm} p{2.5cm} p{4cm}}
\toprule
\textbf{Fase} & \textbf{Operación} & \textbf{Tiempo} & \textbf{Observación} \\
\midrule
FASE 0   & Extracción y segmentación PDF   & 2--3 min   & Depende del tamaño y calidad del PDF \\
FASE 0.5 & Indexación FAISS + BM25         & 3--5 min   & Vectorización text-embedding-004 \\
FASE 1.5 & GraphRAG: extracción tripletas  & 15--20 min & 1 llamada LLM adicional por sección \\
FASE 2   & Análisis multi-agente (41 secc.)& 40--50 min & 3 agentes en paralelo por sección \\
FASE 3   & Consolidación y deduplicación   & 1--2 min   & Postprocesamiento en memoria \\
FASE 4   & Generación reporte PDF          & 1--2 min   & FPDF2, sin llamadas LLM adicionales \\
\midrule
\textbf{Total} & & $\approx \mathbf{75}$\textbf{ min} & \\
\bottomrule
\end{tabular}
\end{table}

La fase de mayor costo es el análisis multi-agente (FASE 2), que representa
aproximadamente el 60\% del tiempo total y escala linealmente con el número
de secciones del contrato.

\section{Métricas Secundarias}

\subsection{Definición, método y resultado}

\begin{table}[H]
\centering
\caption{Métricas secundarias de validación del MVP}
\label{tab:metricas_secundarias}
\begin{tabular}{p{3cm} p{3.2cm} p{2.5cm} p{3.5cm}}
\toprule
\textbf{Métrica} & \textbf{Objetivo SMART} & \textbf{Resultado} & \textbf{Estado} \\
\midrule
Cobertura secciones &
    $100\%$ sin error &
    $100\%$ (41/41) &
    \checkmark\ ALCANZADO \\
Cobertura cláusulas &
    $> 95\%$ analizadas &
    $100\%$ (448/448) &
    \checkmark\ SUPERADO \\
Precisión semántica &
    $\geq 85\%$ hallazgos válidos &
    $\approx 87\%$ (muestra manual) &
    \checkmark\ ALCANZADO \\
Recall estimado &
    $\geq 90\%$ &
    $\approx 94\%$ (estimación) &
    \checkmark\ ALCANZADO \\
Trazabilidad &
    $100\%$ con ubicación + evidencia &
    $100\%$ (119/119) &
    \checkmark\ ALCANZADO \\
Total hallazgos &
    $> 30$ inconsistencias &
    119 (Gemini 2.5 Pro) &
    \checkmark\ SUPERADO ($\times 3.9$) \\
\bottomrule
\end{tabular}
\end{table}

\subsection{Cobertura de secciones y cláusulas}

La cobertura mide qué proporción de las unidades documentales del contrato
fueron procesadas por el pipeline sin error de ejecución:

\[
\text{Cobertura}_{\text{secciones}} = \frac{41}{41} = 100\%
\qquad
\text{Cobertura}_{\text{cláusulas}} = \frac{448}{448} = 100\%
\]

El 100\% de cobertura sobre un contrato con tablas anidadas, anexos técnicos
y numeración mixta valida que el módulo de segmentación es robusto ante la
variabilidad estructural de los contratos de ProInversión.

\subsection{Precisión semántica: protocolo y resultado}

La precisión cuantifica qué proporción de los hallazgos generados
corresponden a inconsistencias reales (verdaderos positivos):

\[
\text{Precisión} = \frac{TP}{TP + FP}
\]

\textbf{Protocolo de medición:}
\begin{enumerate}
    \item Muestra aleatoria estratificada de 30 hallazgos de severidad ALTA
    (estrato de mayor riesgo institucional).
    \item Dos especialistas legales revisaron cada hallazgo de forma
    independiente, sin ver las evaluaciones del otro.
    \item Clasificación: TP (inconsistencia real y bien descrita), FP
    (sin sustento o irrelevante) o \textit{discutible}.
    \item Casos discutibles resueltos por consenso en sesión conjunta.
\end{enumerate}

\textbf{Resultado:} 26 TP, 4 FP $\Rightarrow$ precisión $= 26/30 \approx 87\%$.

Los 4 falsos positivos correspondieron a redacciones ambiguas que el modelo
interpretó como contradictorias pero que los especialistas consideraron
técnicamente correctas en contexto sectorial específico (ingeniería
de saneamiento).

\subsection{Recall estimado}

Dado que no existe un \textit{gold standard} con todas las inconsistencias
del Contrato A anotadas manualmente, el recall no puede medirse
con exactitud. Se realizó la siguiente estimación conservadora:

\begin{itemize}
    \item Durante la revisión de la muestra, los especialistas identificaron
    $\approx 8$ inconsistencias adicionales no detectadas por el sistema
    (falsos negativos), principalmente omisiones de obligaciones implícitas
    que requieren conocimiento sectorial muy específico.
    \item Estimación:
    $\displaystyle\text{Recall} \approx \frac{119}{119 + 8} \approx 94\%$
\end{itemize}

\subsection{Trazabilidad}

Se considera trazable un hallazgo que incluye los cuatro campos siguientes
con valor no vacío:

\begin{itemize}
    \item \textbf{Ubicación exacta:} capítulo, sección y número de cláusula.
    \item \textbf{Evidencia textual:} fragmento literal del contrato.
    \item \textbf{Referencia normativa} (cuando aplica): norma recuperada
    por el módulo RAG.
    \item \textbf{Explicación del razonamiento:} descripción del tipo de
    inconsistencia y por qué es relevante institucionalmente.
\end{itemize}

Los 119 hallazgos del análisis con Gemini~2.5~Pro cumplieron los cuatro
campos, logrando un 100\% de trazabilidad y cumpliendo el requisito
institucional de verificabilidad independiente.

\section{Alineación con Objetivos SMART del Proyecto}

\begin{table}[H]
\centering
\caption{Trazabilidad entre métricas de validación y objetivos SMART (Capítulo 3)}
\label{tab:smart_trazabilidad}
\begin{tabular}{p{4cm} p{3cm} p{2.5cm} p{3cm}}
\toprule
\textbf{Objetivo SMART} & \textbf{Métrica} & \textbf{Resultado} & \textbf{Estado} \\
\midrule
Procesar contrato en $\approx 1$ h &
    Tiempo de screening &
    $\approx 75$ min &
    \checkmark\ SUPERADO \\
Detectar $> 30$ inconsistencias &
    Total hallazgos &
    119 &
    \checkmark\ SUPERADO \\
Precisión $\geq 85\%$ &
    Precisión semántica &
    $\approx 87\%$ &
    \checkmark\ ALCANZADO \\
Recall $\geq 90\%$ &
    Recall estimado &
    $\approx 94\%$ &
    \checkmark\ ALCANZADO \\
Cobertura $100\%$ secciones &
    Cobertura secciones &
    $100\%$ (41/41) &
    \checkmark\ ALCANZADO \\
Trazabilidad $100\%$ &
    Campos trazabilidad &
    $100\%$ (119/119) &
    \checkmark\ ALCANZADO \\
\bottomrule
\end{tabular}
\end{table}

Todos los objetivos SMART fueron alcanzados o superados. El resultado más
destacado es el número de hallazgos (119 vs.\ objetivo $>30$, un factor
$\times 3.9$), que refleja la alta sensibilidad del análisis multi-agente
sobre documentos de esta extensión y complejidad.

\section{Comparación con el Proceso Manual (Baseline)}

\begin{table}[H]
\centering
\caption{ContractIA MVP vs.\ proceso manual de revisión contractual}
\label{tab:comparacion_baseline}
\begin{tabular}{p{3.5cm} p{4cm} p{5cm}}
\toprule
\textbf{Dimensión} & \textbf{Proceso Manual} & \textbf{ContractIA MVP} \\
\midrule
Tiempo total &
    $\approx 320$ h (equipo de 5, 8 semanas) &
    $\approx 75$ min (reducción del $99.6\%$) \\
Cobertura &
    $\approx 60\%$ (muestreo selectivo) &
    $100\%$ secciones y cláusulas \\
Inconsistencias detectadas &
    $\approx 25$ (estimación de especialistas) &
    119 ($\times 4.8$ más hallazgos) \\
Reproducibilidad &
    Baja: depende del especialista &
    Alta: temperatura 0.1; resultados estables \\
Trazabilidad &
    Manual: notas dispersas en Excel &
    Automática: ubicación + evidencia + norma \\
Costo estimado / contrato &
    $\approx \$7{,}200$ (horas-hombre) &
    $\approx \$50$--$\$100$ (tokens + GCP) \\
\bottomrule
\end{tabular}
\end{table}

La reducción de costo por contrato ($> 98\%$) y la multiplicación de
hallazgos detectados ($\times 4.8$) confirman que el MVP genera valor
operativo significativo como herramienta de \textit{screening} previo a
la revisión especializada, sin pretender reemplazar el juicio experto en
la toma de decisiones finales.
